\section{The Jolie Engine and Interpreter}

In this section, the internals of the Jolie engine will be introduced. This
should give the reader the necessary knowledge to understand the changes made
to support a module and package system for Jolie.

At the core of the Jolie engine is the interpreter. Each interpreter is
responsible for parsing and executing a Jolie program. A single engine may run
several instances of the interpreter, this is most commonly the case when
embedding several other Jolie services.

A simplified view of the Jolie interpreter's pipeline can be seen in Figure
\ref{fig:jolie_pipeline}.

\begin{figure}[H]
\centering
\includegraphics[width=0.8\textwidth]{background/pipeline.eps}
\caption{A simplified view of the Jolie interpreter pipeline}
\label{fig:jolie_pipeline}
\end{figure}

The first phase of the interpreter is parsing the intent. In this phase the
interpreter will determine its purpose and, if any, modifiers to its ordinary
behavior.

The Jolie engine can be invoked from the command-line, the command-line is the
source of the intent which starts the first interpreter. The syntax for the
command line is (roughly) as follows: \txtl{jolie [commands and options]
    <program> [program arguments]}.

Options change the overall behavior of the engine, all interpreter instances
share these. These change the behavior of the interpreter.

Commands and program arguments, only belong to the interpreter that they were
originally passed. These determine what the interpreter will do.

The intent parsing phase is responsible for locating and retrieving the program
used.

A URI matching the input program is given to the program parser, the second
phase. The parser is responsible for creating the Abstract Syntax Tree (AST)
    which represents the input program. The parser will produce only a single
    root node, namely the \txtl{Program} node. Includes are thus implemented by
    essentially inserting the contents of the include target. A quirk for Jolie
    includes, is that relative file includes are not resolved relative to the
    current file, but rather to the current working directory, i.e., where the
    engine was started.

The third phase traverses the AST to make sure that it is semantically valid.
This weeds out programs that are syntactically correct, but not semantically
valid. The checks performed in this phase are somewhat limited, given the
dynamic nature of the language.

Given a semantically valid AST the interpreter is ready to build the
interpretation tree. The interpretation tree contains new nodes, which are
runnable.

Once the interpretation tree is built, we're ready to execute the actual code
(which lives in the interpretation tree). The \txtl{CommCore} component is
responsible for all communication, we'll discuss this briefly when relevant.
Following its initialization, the user code of Jolie will start executing. The
\joliel{init} procedure is executed first followed by the \joliel{main}
procedure.
